\begin{table}[!htbp]
\centering
\caption{Multinomial Decomposition of the Binary Estimate, 100,000-Observation Draw}
\label{tab:multinomial}
\resizebox{\ifdim\width>\textwidth \textwidth\else \width\fi}{!}{

\begin{tabular}{lccc}
\toprule
  & Effect (pp per SD) & Share of effect & Share of population\\
\midrule
No high school & -5.48 & --- & ---\\
\midrule
High school & 1.62 & 29.6 & 31.2\\
Some college & 3.12 & 57.0 & 55.5\\
Graduate degree & 0.73 & 13.4 & 13.3\\
\midrule
Implied high school or above & 5.48 & 100.0 & 100.0\\
Binary high school or above & 5.43 & --- & ---\\
\bottomrule
\end{tabular}
}
\begin{minipage}{\textwidth}
\footnotesize
\vspace{4pt}
\textit{Notes:} Two-stage residual inclusion with a four-category outcome. Multinomial coefficients across log-odds equations are not comparable to a binary model's, so both models are converted to average marginal effects on predicted probabilities and rescaled to a one-standard-deviation increase in exposure. The four marginal effects sum to zero by construction. ``Share of effect'' is each completing category's share of the total movement out of no-high-school; ``share of population'' is its share of the completing population in the same draw. Standard errors are not reported: the bootstrap behind the estimate uses 99 metropolitan-area resamples, which is too coarse to report.
\end{minipage}
\end{table}

